%% leaf -- signed leaf cones: the representation and the order, the clause
%% form of every literal, least closures with one seed per atom, duality,
%% extensions first, and collision return. Long proofs are in appendix-leaf.

\section{Signed Leaf Cones}\label{sec:lf}

This section represents trees, and the solutions of systems, by sets of
propositional atoms. In these coordinates every system over $\BFT$, of any
arity and any variance, becomes a set of Horn clauses: implications with one
premise, unit clauses (facts and forbidden roots), and one fixed family of binary
exclusions. The representation has two consequences on which the rest of the
paper rests. Least solutions are computed by closing a set of seeds, and every
atom of a closure has a single seed (\cref{sec:lf-least}). Every pure move
has a dual with both signs flipped, so two derivations that collide can be
merged into one (\cref{sec:lf-duality}).

We use the notation of \cref{sec:bg-trees,sec:bg-constraints}, in
particular the directions $\Sigma=\{0,\dots,n-1\}$, the variance $c$, the
polarity $\parity(w)$, the order $\le_c$, the normalizations $\nu_s$, the dual
$t^\partial$ and the full constructor tree $\Omega$, and we put
$\Sigma^+=\Sigma^*\setminus\{\eps\}$.
In this section and the next, all trees are trees over $\BFT$, and
assignments range over arbitrary trees; regular and finite trees are treated
in \cref{sec:tc}. Proofs omitted here are in \cref{app:lf}.

The positions of a tree $t$ labeled $\Bot$ or $\Top$ are its \emph{leaves}. We
use the \emph{chain property}: for every tree $t$ and every word $w$, at most
one prefix of $w$ is a leaf of $t$, and if none is, then $w$ is a position of
$t$ labeled $f$. Indeed, by induction on its prefixes, a prefix of $w$ is a
position as long as no shorter prefix is a leaf, and a position that is not a
leaf has all its children.

\subsection{Leaf Cones}\label{sec:lf-cones}

\begin{definition}[Signed leaf cones]\label{def:lf-cones}
For a tree $T$, the \emph{top-leaf cone} $A_T$ is the set of words $w$ such
that some prefix of $w$ is a position of $T$ labeled $\Top$. For a tree $t$
and a sign $s\in\{0,1\}$, the \emph{leaf cone of sign $s$} of $t$ is
\[
 L^s_t=A_{\nu_st}.
\]
We also write $L^+_t=L^0_t$ and $L^-_t=L^1_t$, and call $(L^0_t,L^1_t)$ the
\emph{leaf pair} of $t$.
\end{definition}

The cones have a direct description. The \emph{sign} of a leaf $v$ of $t$ is
$b\oplus\parity(v)$, where $b=0$ if $t(v)=\Top$ and $b=1$ if $t(v)=\Bot$.
The tree $\nu_st$ carries the label $\Top$ at a leaf $v$ iff $v$ has sign $s$.
By the chain property, therefore,
\begin{equation}\label{eq:lf-direct}
 w\in L^s_t\quad\Longleftrightarrow\quad
 \text{the path $w$ meets a leaf of $t$ of sign $s$.}
\end{equation}
Since $\nu_1t=(\nu_0t)^\partial$ (\cref{lem:bg-normalization}(1)), the cone
$L^1_t$ is the bottom-leaf cone of $\nu_0t$: the two cones belong to the same normalized tree. The variance
enters only through the sign of each leaf, and the parity that determines this
sign accumulates along the path and stops at the leaf. There are two cones per
tree, whatever the arity, not one per assignment of signs to the children.
In terms of the signed correspondence (\cref{lem:bg-signed}),
$L^s_{\rho(v)}$ is the top-leaf cone of the covariant tree $\rho^\pm(v^s)$.

\begin{theorem}[Leaf representation]\label{thm:lf-representation}
\leavevmode
\begin{enumerate}[leftmargin=*]
\item For every tree $t$, the sets $L^0_t$ and $L^1_t$ are disjoint, and each
  is \emph{extension closed}: $w\in L^s_t$ implies $wi\in L^s_t$ for every
  letter $i$. Conversely, every pair $(L^0,L^1)$ of disjoint extension-closed
  sets of words is the leaf pair of exactly one tree.
\item For all trees $t$ and $u$,
  \[
   t\le_c u\quad\Longleftrightarrow\quad
   L^0_t\subseteq L^0_u\ \text{ and }\ L^1_u\subseteq L^1_t .
  \]
\item $L^0_\Top=L^1_\Bot=\Sigma^*$, $L^1_\Top=L^0_\Bot=\varnothing$, and
  $L^s_{t^\partial}=L^{1-s}_t$. If $t=f(t_i)_i$, then neither cone of $t$
  contains $\eps$, and for every letter $i$, word $w$, and sign $s$,
  \[
   iw\in L^s_t\quad\Longleftrightarrow\quad w\in L^{s\oplus c(i)}_{t_i}.
  \]
  In particular, $\eps\in L^0_t$ iff $t=\Top$, and $\eps\in L^1_t$ iff
  $t=\Bot$.
\item The positions of $t$ labeled $f$ are exactly the words outside
  $L^0_t\cup L^1_t$. Hence $t$ is finite iff
  $\Sigma^*\setminus(L^0_t\cup L^1_t)$ is finite, and $t$ is regular iff
  $L^0_t$ and $L^1_t$ are regular languages.
\end{enumerate}
\end{theorem}

By (2), the order adds leaves of sign $0$ and removes leaves of sign $1$ as
it goes up: after normalization, larger trees have more top leaves and fewer
bottom leaves. The order is therefore not the inclusion of both cones.
Inclusion of both cones is an \emph{information order}, whose least element
is the tree without leaves, $\Omega$, with the leaf pair
$(\varnothing,\varnothing)$. The least solutions of \cref{sec:lf-least} are
least in the information order.

\begin{example}[Leaf cones of arrows]\label{ex:lf-arrow}
For the arrow, $n=2$ and $f(a_0,a_1)=a_0\to a_1$ with $c(0)=1$ and $c(1)=0$, so
$\parity(w)$ is the parity of the number of letters $0$ in $w$. The leaf pairs
of $\Top$, $\Bot$ and $\Omega$ are $(\Sigma^*,\varnothing)$,
$(\varnothing,\Sigma^*)$ and $(\varnothing,\varnothing)$. In $\Top\to\Bot$, the
leaf $0$ is labeled $\Top$ and has polarity $1$, and the leaf $1$ is labeled
$\Bot$ and has polarity $0$; both have sign $1$, so the leaf pair is
$(\varnothing,\Sigma^+)$. Likewise $\Bot\to\Top$ has the leaf pair
$(\Sigma^+,\varnothing)$, and $\Bot\to\Omega$ has $(0\Sigma^*,\varnothing)$. By
\cref{thm:lf-representation}(3), both cones of every tree $f(a_0,a_1)$ lie in
$\Sigma^+$, so (2) gives $\Top\to\Bot\le_c f(a_0,a_1)\le_c\Bot\to\Top$: these
are the least and the greatest arrow types.
\end{example}

\Cref{fig:lf-cones} follows three arrow types up to the greatest one,
$\Bot\to\Top$. The sign of each leaf is read off its label and its path, and
each step up either removes a leaf of sign~$1$ or adds a leaf of sign~$0$, even
where a subtree at a contravariant position becomes smaller.

\begin{figure}[tbp]
\centering
\begin{tikzpicture}[>=Stealth,font=\small,
    every node/.style={inner sep=1.5pt},
    tree/.style={level distance=8mm,
      level 1/.style={sibling distance=16mm},
      level 2/.style={sibling distance=9mm}},
    e0/.style={left,font=\scriptsize,inner sep=1pt,pos=0.45},
    e1/.style={right,font=\scriptsize,inner sep=1pt,pos=0.45},
    s0/.style={draw,rounded corners=1pt,fill=black!12,minimum size=4.6mm},
    s1/.style={draw,rounded corners=1pt,minimum size=4.6mm},
    sg/.style={below=0.5mm,font=\scriptsize},
    lb/.style={font=\footnotesize,align=left,anchor=north west}]
  % t1 = (Bot -> Bot) -> Top
  \begin{scope}[tree]
    \node at (1.3,0) {$f$}
      child {node {$f$}
        child {node[s1] (l00) {$\Bot$} edge from parent node[e0]{0}}
        child {node[s0] (l01) {$\Bot$} edge from parent node[e1]{1}}
        edge from parent node[e0]{0}}
      child {node[s0] (l1) {$\Top$} edge from parent node[e1]{1}};
  \end{scope}
  \node[sg] at (l00.south) {$1$};
  \node[sg] at (l01.south) {$0$};
  \node[sg] at (l1.south) {$0$};
  \node[font=\footnotesize,anchor=south] at (1.3,0.15) {$t_1=(\Bot\to\Bot)\to\Top$};
  \node[lb] at (0.0,-2.15) {$L^0_{t_1}=01\Sigma^*\cup1\Sigma^*$\\ $L^1_{t_1}=00\Sigma^*$};
  % t2 = (Omega -> Bot) -> Top
  \begin{scope}[tree]
    \node at (6.2,0) {$f$}
      child {node {$f$}
        child {node (o00) {$\Omega$} edge from parent node[e0]{0}}
        child {node[s0] (m01) {$\Bot$} edge from parent node[e1]{1}}
        edge from parent node[e0]{0}}
      child {node[s0] (m1) {$\Top$} edge from parent node[e1]{1}};
  \end{scope}
  \node[sg] at (m01.south) {$0$};
  \node[sg] at (m1.south) {$0$};
  \node[font=\footnotesize,anchor=south] at (6.2,0.15) {$t_2=(\Omega\to\Bot)\to\Top$};
  \node[lb] at (4.9,-2.15) {$L^0_{t_2}=01\Sigma^*\cup1\Sigma^*$\\ $L^1_{t_2}=\varnothing$};
  % t3 = Bot -> Top
  \begin{scope}[tree]
    \node at (10.9,0) {$f$}
      child {node[s0] (n0) {$\Bot$} edge from parent node[e0]{0}}
      child {node[s0] (n1) {$\Top$} edge from parent node[e1]{1}};
  \end{scope}
  \node[sg] at (n0.south) {$0$};
  \node[sg] at (n1.south) {$0$};
  \node[font=\footnotesize,anchor=south] at (10.9,0.15) {$t_3=\Bot\to\Top$};
  \node[lb] at (9.9,-2.15) {$L^0_{t_3}=\Sigma^+$\\ $L^1_{t_3}=\varnothing$};
  % the order
  \node at (3.75,-0.8) {$\le_c$};
  \node[font=\scriptsize,align=center] at (3.75,-1.35) {sign-$1$ leaf\\ $00$ removed};
  \node at (8.65,-0.8) {$\le_c$};
  \node[font=\scriptsize,align=center] at (8.65,-1.35) {sign-$0$ leaf\\ $0$ added};
\end{tikzpicture}
\caption{Signed leaf cones of three arrow types, $c(0)=1$ and $c(1)=0$. The
number below a leaf is its sign: $0$ for $\Top$ and $1$ for $\Bot$, plus the
number of letters $0$ on its path, modulo two. Leaves of sign $0$ are shaded;
the leaf $\Bot$ at $01$ has sign $0$, and so does the leaf $\Bot$ at $0$ in
$t_3$. Going up in the order removes leaves of sign $1$ and adds leaves of sign
$0$ (\cref{thm:lf-representation}(2)): $t_1\le_ct_2$ shrinks $L^1$, and
$t_2\le_ct_3$ grows $L^0$ by replacing the subtree $\Omega\to\Bot$ at the
contravariant position $0$ with the smaller tree $\Bot$.}
\Description{Three trees side by side, separated by the symbol less or equal
for the variance c. Left: the arrow type from bottom to bottom, to top; its
leaves are bottom at 00 with sign 1, bottom at 01 with sign 0, and top at 1
with sign 0. Middle: the same tree with Omega in place of the leaf at 00; its
leaves bottom at 01 and top at 1 both have sign 0. Right: bottom to top, whose
two leaves have sign 0. Under each tree its two leaf cones are listed: zero one
Sigma star union one Sigma star and zero zero Sigma star; zero one Sigma star
union one Sigma star and the empty set; Sigma plus and the empty set. Leaves of
sign 0 are shaded. Between the first two trees: sign 1 leaf 00 removed; between
the last two: sign 0 leaf 0 added.}
\label{fig:lf-cones}
\end{figure}


\subsection{Atoms and Clauses}\label{sec:lf-clauses}

Let $V$ be a finite set of variables, and recall the signed variables $v^s$,
their set $V^\pm$, and the flip $\overline{v^s}=v^{1-s}$
(\cref{sec:bg-constraints}); we put $\sgn(v^s)=s$. An \emph{atom} is a pair
$(p,w)\in\Conf_V=V^\pm\times\Sigma^*$, and its \emph{opposite} is
$\overline{(p,w)}=(\bar p,w)$. An assignment $\rho$ of trees to $V$
determines the set of atoms
\[
 S_\rho=\{(v^s,w):w\in L^s_{\rho(v)}\},
\]
and an atom \emph{holds} in $\rho$ if it lies in $S_\rho$. An atom of sign
$0$ is \emph{positive}, and an atom of sign $1$ is \emph{negative}. Opposite atoms are
incompatible but not complementary: in $\Omega$ neither holds. For
$U\subseteq V$ and $S\subseteq\Conf_V$ we put $S|_U=S\cap\Conf_U$. We read
the word of an atom as a stack whose top is its first letter.

A set $S$ of atoms is \emph{extension closed} if $(p,w)\in S$ implies
$(p,wi)\in S$ for every letter $i$, and an \emph{exclusion pair} of $S$ is a
pair $a,\bar a$ of opposite atoms that both lie in $S$. By
\cref{thm:lf-representation}(1), applied to each variable, $\rho\mapsto
S_\rho$ is a bijection from the assignments of trees to $V$ onto the
extension-closed subsets of $\Conf_V$ without exclusion pairs.

\begin{table}[t]
\caption{The clause form of a literal, for all words $w$ and all letters
$i\in\Sigma$; here $\bar c(i)=1\oplus c(i)$. Every system also has the
extensions $(p,w)\to(p,wi)$ for all $p\in V^\pm$, $w\in\Sigma^*$ and
$i\in\Sigma$.}
\label{tab:lf-clauses}
\centering
\begin{tabular}{@{}lll@{}}
\toprule
Literal & Pure moves & Root unit\\
\midrule
$u\le v$ & copies: $(u^0,w)\to(v^0,w)$ and $(v^1,w)\to(u^1,w)$ & none\\[2pt]
$f(a_i)_i\le u$ & $\mathrm{push}_i$: $(a_i^{c(i)},w)\to(u^0,iw)$
 & forbidden root $(u^1,\eps)$\\
 & $\mathrm{pop}_i$: $(u^1,iw)\to(a_i^{\bar c(i)},w)$ & \\[2pt]
$u\le f(a_i)_i$ & $\mathrm{pop}_i$: $(u^0,iw)\to(a_i^{c(i)},w)$
 & forbidden root $(u^0,\eps)$\\
 & $\mathrm{push}_i$: $(a_i^{\bar c(i)},w)\to(u^1,iw)$ & \\[2pt]
$u=\Bot$ & none & fact $(u^1,\eps)$\\
$\Top\le u$ & none & fact $(u^0,\eps)$\\
\bottomrule
\end{tabular}
\end{table}

The literals of a system $\chi$ over $V$ determine the \emph{pure moves},
\emph{facts} and \emph{forbidden roots} listed in
\cref{tab:lf-clauses}; facts and forbidden roots are the unit clauses. A push or a pop of direction $i$ changes the sign by
$c(i)$, and a copy keeps it. The \emph{rules} of $\chi$ are its pure moves and
the \emph{extensions} $(p,w)\to(p,wi)$. We write $\Fact_\chi$ and $\Neg_\chi$
for the sets of facts and forbidden roots. Pure moves read and write
only the top of the stack: $\mathrm{push}_i$ puts the letter $i$ on top,
$\mathrm{pop}_i$ removes it, and a copy leaves the stack unchanged. They form a
pushdown system with control states $V^\pm$ and stack alphabet $\Sigma$,
which \cref{sec:sm} summarizes by finite automata. Extensions, in contrast,
add a letter at the bottom of the stack.

The table is the signed system $\chi^\pm$ of \cref{tab:bg-signed} read in
top-leaf cones. Every copy literal $p\le q$ of $\chi^\pm$ gives the pure moves
$(p,w)\to(q,w)$, every lower constructor literal $f(p_i)_i\le q$ gives
$(p_i,w)\to(q,iw)$, and every upper constructor literal $q\le f(p_i)_i$ gives
$(q,iw)\to(p_i,w)$. The forbidden roots are the atoms $(q,\eps)$ for the roots
$q$ of the upper constructor literals of $\chi^\pm$, and the facts
are the atoms $(q,\eps)$ for its literals $\Top\le q$.

\begin{definition}[Models]\label{def:lf-model}
A set $S\subseteq\Conf_V$ is \emph{consistent} for $\chi$ if it contains no
forbidden root of $\chi$ and no exclusion pair. A \emph{model} of $\chi$ is a
consistent set of atoms that contains $\Fact_\chi$ and is closed under all
rules of $\chi$. We write $\Mod(\chi)$ for the set of models of $\chi$.
\end{definition}

\begin{lemma}[Clause form]\label{lem:lf-clauses}
An assignment $\rho$ satisfies a literal iff $S_\rho$ is closed under the
pure moves of the literal, contains its facts, and contains none of
its forbidden roots. Consequently, $\rho\mapsto S_\rho$ is a bijection from
the solutions of $\chi$ onto $\Mod(\chi)$.
\end{lemma}

The proof (\cref{app:lf}) substitutes the identities of
\cref{thm:lf-representation}(3) into the inclusions of
\cref{thm:lf-representation}(2). Every literal thus contributes implications
with one premise and units only; the only all-negative binary clauses are the
exclusions $\neg(a\wedge\bar a)$, which are the same for all systems, and no
clause is a positive disjunction. A constructor literal contributes $2n$
schemes of pure moves and one forbidden root, and a copy literal two schemes.
For $n=0$, the literal $f()\le u$ says $u\ne\Bot$ and $u\le f()$ says
$u\ne\Top$.

We write $a\to_\chi b$ if a pure move of $\chi$ leads from $a$ to $b$.

\begin{lemma}[Pure moves]\label{lem:lf-moves}
Let $(p,u)\to_\chi(q,t)$.
\begin{enumerate}[leftmargin=*]
\item For every word $\sigma$, $(p,u\sigma)\to_\chi(q,t\sigma)$, by the same
  scheme of \cref{tab:lf-clauses}.
\item $\sgn(q)\oplus\parity(t)=\sgn(p)\oplus\parity(u)$.
\end{enumerate}
\end{lemma}

Item (1) holds because each scheme is stated for all words $w$, and (2) is a
direct check of the six schemes (\cref{app:lf}). Both extend to sequences
of pure moves. By (2), along such a sequence the sign of the current signed
variable is determined by the first atom and the parity of the current word;
we call (2) the \emph{sign invariant}.

\begin{example}[The system $x\le x\to x$]\label{ex:lf-clauses}
Let $\varphi=\{x\le x\to x\}$ for the arrow of \cref{ex:lf-arrow}, that is,
the literal $x\le f(x,x)$. By \cref{tab:lf-clauses}, its pure moves are
\[
\begin{array}{ll}
 \mathrm{pop}_0:\ (x^0,0w)\to(x^1,w), & \mathrm{push}_0:\ (x^0,w)\to(x^1,0w),\\
 \mathrm{pop}_1:\ (x^0,1w)\to(x^0,w), & \mathrm{push}_1:\ (x^1,w)\to(x^1,1w),
\end{array}
\]
it has the forbidden root $(x^0,\eps)$, and it has no facts. A move
in the contravariant direction $0$ changes the sign, and a move in the
covariant direction $1$ keeps it. In each row, the push is the pop read
backwards with both signs reversed; \cref{lem:lf-duality} shows that this
holds for every system. The tree $\Top\to\Bot$, with leaf pair
$(\varnothing,\Sigma^+)$, is a solution: its atoms are the $(x^1,w)$ with
$w\ne\eps$, and this set is closed under $\mathrm{push}_1$ and under
extensions.
\end{example}

\subsection{Least Closures}\label{sec:lf-least}

A \emph{derivation} of $\chi$ is a finite sequence of atoms $a_0,\dots,a_k$,
with $k\ge0$, in which each $a_{j+1}$ follows from $a_j$ by a rule of $\chi$;
it is \emph{pure} if all its rules are pure moves. We write
$a\Rightarrow^*_\chi b$ if a derivation leads from $a$ to $b$, and
$a\to^*_\chi b$ if a pure derivation does; we drop the subscript when $\chi$
is clear. For an atom $a$ and a set $A$ of atoms put
\[
 \Reach_\chi(a)=\{b:a\Rightarrow^*_\chi b\},\qquad
 \mu_\chi=\bigcup_{e\in\Fact_\chi}\Reach_\chi(e),\qquad
 H_\chi(A)=\mu_\chi\cup\bigcup_{a\in A}\Reach_\chi(a),
\]
and $\Cl_\chi(a)=H_\chi(\{a\})$. An atom is \emph{live} in $\chi$ if some
model of $\chi$ contains it, and \emph{dead} otherwise; $\Live_\chi$ is the
set of live atoms. Every rule has a single premise, so every atom $b$ of
$H_\chi(A)$ has a \emph{seed}: an atom $x\in\Fact_\chi\cup A$ with
$b\in \Reach_\chi(x)$.

\begin{lemma}[Least closures]\label{lem:lf-closure}\label{lem:lf-least}
Let $\chi$ be a system over $V$ and $A\subseteq\Conf_V$.
\begin{enumerate}[leftmargin=*]
\item $H_\chi(A)$ is the least set of atoms that contains $\Fact_\chi\cup A$
  and is closed under the rules of $\chi$.
\item Some model of $\chi$ contains $A$ iff $H_\chi(A)$ is consistent, and
  then $H_\chi(A)$ is the least model that contains $A$. In particular,
  $\chi$ is satisfiable iff $\mu_\chi$ is consistent, and then $\mu_\chi$ is
  the least model of $\chi$; and an atom $a$ is live iff $\Cl_\chi(a)$ is
  consistent, and then $\Cl_\chi(a)$ is the least model that contains $a$.
\item If $H_\chi(A)$ is not consistent, then there are seeds
  $x,y\in\Fact_\chi\cup A$, possibly equal, such that $\Reach_\chi(x)$ contains a
  forbidden root, or $b\in \Reach_\chi(x)$ and $\bar b\in \Reach_\chi(y)$ for some atom
  $b$.
\item For every atom $(p,u)$ and every word $\alpha$,
  $(p,u\alpha)\in \Reach_\chi(p,u)$; hence
  $\Cl_\chi(p,u\alpha)\subseteq\Cl_\chi(p,u)$, and if $(p,u)$ is live, then
  so is $(p,u\alpha)$.
\item Let $V=X\cup Z$ with $X\cap Z=\varnothing$, and let $S\subseteq\Conf_X$
  be extension closed and without exclusion pairs. Then $S=T|_X$ for some
  model $T$ of $\chi$ iff $H_\chi(S)$ is consistent and $H_\chi(S)|_X=S$; in
  this case $H_\chi(S)$ is the least such model.
\end{enumerate}
\end{lemma}

Item (3) says that every inconsistency has at most two seeds; it bounds
neither the words of the seeds nor the lengths of the derivations. Item (5)
reduces the extension of a given visible assignment to a single closure;
\cref{thm:ch-characterization} and the transfers of \cref{sec:tc} use the
lemma in this form.

\begin{example}[Liveness in $x\le x\to x$]\label{ex:lf-least}
We continue \cref{ex:lf-clauses}. The system $\varphi$ has no facts,
so $\mu_\varphi=\varnothing$ is consistent: $\varphi$ is satisfiable, and its
least solution is $x=\Omega$. Every atom $(x^1,u)$ is live, since the
solution $x=\Bot$ contains it. The atoms $(x^0,1^k)$ are dead: $k$ pops in
direction $1$ lead to the forbidden root $(x^0,\eps)$. The atom $(x^0,0)$ is
dead for a different reason. The pop $(x^0,0)\to(x^1,\eps)$ followed by an
extension derives $(x^1,0)$, the opposite of the seed, so
$\Cl_\varphi((x^0,0))$ contains an exclusion pair. No forbidden root is
derivable from $(x^0,0)$: the only atoms $(x^0,w)$ that it derives are its
extensions $(x^0,0w)$, and no move leads from a negative atom of $x$ to a
positive one. So this inconsistency has a single seed and uses no unit.
Finally, $(x^0,01)$ is live: its least model is the solution
$((\Omega\to\Bot)\to\Bot)\to\Bot$, whose atoms are the $(x^0,w)$ with
$w\in01\Sigma^*$ and the $(x^1,w)$ with $w\in1^+\Sigma^*\cup1^*001\Sigma^*$.
\end{example}

\Cref{fig:lf-moves}(a) draws these derivations with each atom at the height of
its word, so that the moves act as on a stack: pops go down, pushes go up, and
an extension adds a letter at the bottom without changing the top.

\begin{figure}[tbp]
\centering
\begin{tikzpicture}[>=Stealth,font=\small,x=1cm,y=1cm,
    every node/.style={inner sep=1.5pt},
    atom/.style={inner sep=2pt},
    seed/.style={atom,draw,thick,rounded corners=1pt},
    hidden/.style={atom,draw,dashed,rounded corners=1pt},
    mv/.style={->},
    ml/.style={font=\scriptsize,inner sep=1pt},
    opp/.style={densely dotted,thick}]
  % height axis
  \draw[black!40] (-0.75,-0.25) -- (-0.75,2.45);
  \foreach \h/\y in {0/0,1/1.1,2/2.2}
    {\draw[black!40] (-0.82,\y) -- (-0.68,\y);
     \node[font=\scriptsize,anchor=east] at (-0.82,\y) {$\h$};}
  \node[font=\scriptsize,anchor=south] at (-0.75,2.45) {$\len w$};
  % (a) phi = x <= x -> x
  \node[font=\footnotesize,anchor=west] at (-0.5,3.0) {(a) $\varphi=x\le x\to x$};
  \node[seed] (s) at (0,1.1) {$(x^0,0)$};
  \node[atom] (p) at (1.15,0) {$(x^1,\eps)$};
  \node[atom] (o) at (2.3,1.1) {$(x^1,0)$};
  \node[atom] (e) at (1.15,2.2) {$(x^0,00)$};
  \draw[mv] (s) -- node[ml,below left]{$\mathrm{pop}_0$} (p);
  \draw[mv] (p) -- node[ml,below right]{ext} (o);
  \draw[mv,dashed] (s) -- node[ml,above left]{ext} (e);
  \draw[mv,dashed] (e) -- node[ml,above right]{$\mathrm{pop}_0$} (o);
  \draw[opp] (s) -- node[ml,above,fill=white]{opposites} (o);
  \node[atom] (r2) at (3.35,2.2) {$(x^0,11)$};
  \node[atom] (r1) at (3.9,1.1) {$(x^0,1)$};
  \node[atom,draw,double,rounded corners=1pt] (r0) at (4.45,0) {$(x^0,\eps)$};
  \draw[mv] (r2) -- node[ml,right]{$\mathrm{pop}_1$} (r1);
  \draw[mv] (r1) -- node[ml,right]{$\mathrm{pop}_1$} (r0);
  \node[ml,anchor=north] at (4.45,-0.25) {forbidden root};
  % (b) psi = x <= z -> z
  \node[font=\footnotesize,anchor=west] at (5.6,3.0) {(b) $\psi=x\le z\to z$};
  \node[seed] (a) at (6.0,1.1) {$(x^0,0)$};
  \node[seed] (b) at (11.6,1.1) {$(x^0,1)$};
  \node[hidden] (z1) at (7.3,0) {$(z^1,\eps)$};
  \node[hidden] (z0) at (9.9,0) {$(z^0,\eps)$};
  \node[atom] (xb) at (8.8,1.1) {$(x^1,1)$};
  \draw[mv,very thick] (a) -- node[ml,below left]{$\mathrm{pop}_0$} (z1);
  \draw[mv] (b) -- node[ml,below right]{$\mathrm{pop}_1$} (z0);
  \draw[mv,very thick] (z1) -- node[ml,below right]{$\mathrm{push}_1$} (xb);
  \draw[opp] (z1) -- node[ml,below=1pt]{opposites} (z0);
  \draw[opp] (xb) -- node[ml,above=1pt]{opposites} (b);
  \node[ml,align=center] at (8.8,2.2) {$\mathrm{push}_1$ is the dual of $\mathrm{pop}_1$};
\end{tikzpicture}
\caption{Pure moves and extensions for the arrow, with each atom drawn at the
height $\len w$ of its word: a pop lowers the stack, a push raises it, and an
extension adds a letter at the bottom. Seeds are framed, hidden atoms dashed.
(a) In $x\le x\to x$ (\cref{ex:lf-clauses,ex:lf-least}), the seed $(x^0,0)$
derives its own opposite $(x^1,0)$, by a pop and an extension or, dashed, with
the extension first (\cref{lem:lf-ext-first}): a single-seed inconsistency
without any unit. The atom $(x^0,11)$ is dead for a different reason: it pops
down to the forbidden root $(x^0,\eps)$. (b) In $x\le z\to z$
(\cref{ex:ch-arrow}), the seeds $(x^0,0)$ and $(x^0,1)$ derive the opposite
atoms $(z^1,\eps)$ and $(z^0,\eps)$ of the hidden $z$. Reversing the second
branch with flipped signs (\cref{lem:lf-duality}) continues the first one (thick):
the single seed $(x^0,0)$ derives $(x^1,1)$, the opposite of the second seed,
which is collision return (\cref{lem:lf-collision}).}
\Description{Two panels of atoms placed at heights 0, 1 and 2, the length of
their words. Panel a: the framed seed x zero at word 0 pops to x one at the
empty word, which extends to x one at word 0, drawn at the height of the seed
and joined to it by a dotted line labeled opposites; a dashed route first extends the seed to x
zero at 00 and then pops to the same atom. To the right, x zero at 11 pops
twice to the doubly framed forbidden root x zero at the empty word. Panel b:
the framed seeds x zero at 0 and x zero at 1 pop to the dashed hidden atoms z
one and z zero at the empty word, joined by a dotted line labeled opposites;
a thick push in direction 1 leads from z one to x one at word 1, which is
joined by a dotted line to the second seed.}
\label{fig:lf-moves}
\end{figure}


\subsection{Duality, Extensions First, and Collision Return}\label{sec:lf-duality}

\begin{lemma}[Duality]\label{lem:lf-duality}
If $a\to_\chi b$, then $\bar b\to_\chi\bar a$. Consequently, $a\to^*_\chi b$
implies $\bar b\to^*_\chi\bar a$.
\end{lemma}

\begin{proof}
By \cref{tab:lf-clauses}: the two copies of a literal $u\le v$ are each
other's duals, and so are $\mathrm{push}_i$ and $\mathrm{pop}_i$ of the same
constructor literal, for each direction $i$ separately; for instance, the
dual of $(a_i^{c(i)},w)\to(u^0,iw)$ is $(u^1,iw)\to(a_i^{\bar c(i)},w)$.
This holds also when variables are repeated among $u,a_0,\dots,a_{n-1}$,
since the pairing is by scheme. For derivations, reverse the sequence and
dualize each move.
\end{proof}

Extensions have no duals: the reverse of an extension would delete the
bottom letter of a stack, and no rule does that. The next lemma moves all
extensions of a derivation to its beginning, so that duality applies to the
rest.

\begin{lemma}[Extensions first]\label{lem:lf-ext-first}\label{lem:lf-extensions}
For all atoms $(p,u)$ and $b$, $b\in \Reach_\chi(p,u)$ iff
$(p,u\alpha)\to^*_\chi b$ for some word $\alpha$.
\end{lemma}

The proof (\cref{app:lf}) inducts on the derivation: an extension at the
end appends a letter to the bottom of every stack of the pure derivation
obtained so far, which \cref{lem:lf-moves}(1) allows.

\begin{lemma}[Collision return]\label{lem:lf-collision}
Let $a$ and $(q,v)$ be atoms. If $b\in \Reach_\chi(a)$ and $\bar b\in \Reach_\chi(q,v)$
for some atom $b$, then $(\bar q,v\delta)\in \Reach_\chi(a)$ for some word
$\delta$. If moreover $a=(p,u)$, then
$(p,u\alpha)\to^*_\chi(\bar q,v\delta)$ for some words $\alpha$ and $\delta$.
\end{lemma}

\begin{proof}
By \cref{lem:lf-ext-first}, $(q,v\delta)\to^*\bar b$ for some $\delta$, and by
\cref{lem:lf-duality}, $b\to^*(\bar q,v\delta)$. Appending this pure
derivation to a derivation from $a$ to $b$ gives
$(\bar q,v\delta)\in \Reach_\chi(a)$. For the second claim apply
\cref{lem:lf-ext-first} to this membership.
\end{proof}

Thus two derivations that end in an exclusion pair, possibly at a hidden
variable, merge into one derivation from the first seed to the opposite of
an extension of the second seed (\cref{fig:lf-moves}(b)): the first
derivation continues along the second one, read backwards with flipped signs.
The extension is necessary: in general the
merged derivation does not reach the opposite of the second seed itself
(\cref{ex:ch-arrow}). The first consequence concerns a single seed.

\begin{corollary}[Deadness]\label{cor:lf-dead}
Let $\chi$ be satisfiable and let $(p,u)$ be an atom. Then $(p,u)$ is dead
in $\chi$ iff at least one of the following holds:
\begin{enumerate}[label=(\alph*),leftmargin=*]
\item $\Reach_\chi(p,u)$ contains a forbidden root;
\item $(\bar p,u\delta)\in\mu_\chi$ for some word $\delta$;
\item $(\bar p,u\delta)\in \Reach_\chi(p,u)$ for some word $\delta$.
\end{enumerate}
Moreover, (a) holds iff $(p,u\alpha)\to^*_\chi e$ for some word $\alpha$ and
some $e\in\Neg_\chi$, and (c) holds iff
$(p,u\alpha)\to^*_\chi(\bar p,u\delta)$ for some words $\alpha$ and $\delta$.
\end{corollary}

Causes (a) and (b) make the atom dead by support (\cref{def:sm-dead}).
Cause (c), a \emph{self-collision}, may occur without such support: in
\cref{ex:lf-least}, $(x^0,0)$ is dead by $(x^0,00)\to(x^1,0)$, which is (c)
with $\alpha=0$ and $\delta=\eps$. By the sign invariant, (c) requires
$\parity(\alpha)\ne\parity(\delta)$, so self-collisions occur only through
contravariant directions; in the covariant case every rule keeps the
sign, and only (a) and (b) remain (\cref{rem:ch-parity}). The causes may
overlap: for unary contravariant $f$ and $\chi=\{x\le f(x)\}$, the atom
$(x^0,\eps)$ is a forbidden root, and also $(x^0,0)\to(x^1,\eps)$. \Cref{sec:sm}
expresses (a)--(c) by finite automata and reduces a self-collision of an
atom that is not already dead by support to a \emph{self event}
(\cref{def:sm-dead}): two cuts $i<j$ of the seed word, the suffix after $j$ a
prefix of the suffix after $i$, such that some signed variable that the seed
reaches by popping the letters before $i$ has its flip reached by popping
the letters before $j$.
